\begin{tikzpicture}[x=1cm,y=1cm,every node/.style={font=\sffamily\fontsize{9.2}{11}\selectfont,text=Ink}]
\node[draw=Classical,fill=Classical,rounded corners=2mm,text=white,text width=2.45cm,align=flush center,minimum height=1.55cm,anchor=west,font=\sffamily\bfseries\fontsize{9.2}{11}\selectfont] (f0) at (0,0.00) {1. Definire l’evento};
\node[draw=Classical!60,fill=Classical!6,rounded corners=1.5mm,line width=.6pt,text width=2.8cm,align=flush center,minimum height=1.55cm,anchor=west] (e00) at (2.95,0.00) {Dalle osservazioni ai massimi annuali\\[1pt]{\fontsize{7.6}{9}\selectfont\color{Muted}m14-massimi}};
\node[draw=Rule,fill=white,rounded corners=1.5mm,line width=.6pt,text width=2.8cm,align=flush center,minimum height=1.55cm,anchor=west] (e01) at (6.10,0.00) {Superamenti e separazione degli episodi\\[1pt]{\fontsize{7.6}{9}\selectfont\color{Muted}m14-episodi}};
\draw[pEdge] (f0.east)--(e00.west);
\node[draw=Classical,fill=Classical,rounded corners=2mm,text=white,text width=2.45cm,align=flush center,minimum height=1.55cm,anchor=west,font=\sffamily\bfseries\fontsize{9.2}{11}\selectfont] (f1) at (0,-1.95) {2. Stimare e ricampionare};
\draw[pEdge] (f0.south)--(f1.north);
\node[draw=Classical!60,fill=Classical!6,rounded corners=1.5mm,line width=.6pt,text width=2.8cm,align=flush center,minimum height=1.55cm,anchor=west] (e10) at (2.95,-1.95) {Adattamento di Gumbel ai massimi\\[1pt]{\fontsize{7.6}{9}\selectfont\color{Muted}m14-gumbel}};
\node[draw=Rule,fill=white,rounded corners=1.5mm,line width=.6pt,text width=2.8cm,align=flush center,minimum height=1.55cm,anchor=west] (e11) at (6.10,-1.95) {Livelli di ritorno e incertezza bootstrap\\[1pt]{\fontsize{7.6}{9}\selectfont\color{Muted}m14-livelli-rari}};
\draw[pEdge] (f1.east)--(e10.west);
\node[draw=Classical,fill=Classical,rounded corners=2mm,text=white,text width=2.45cm,align=flush center,minimum height=1.55cm,anchor=west,font=\sffamily\bfseries\fontsize{9.2}{11}\selectfont] (f2) at (0,-3.90) {3. Interpretare il rischio};
\draw[pEdge] (f1.south)--(f2.north);
\node[draw=Classical!60,fill=Classical!6,rounded corners=1.5mm,line width=.6pt,text width=2.8cm,align=flush center,minimum height=1.55cm,anchor=west] (e20) at (2.95,-3.90) {Sensibilità alla forma della coda GEV\\[1pt]{\fontsize{7.6}{9}\selectfont\color{Muted}m14-gev}};
\node[draw=Rule,fill=white,rounded corners=1.5mm,line width=.6pt,text width=2.8cm,align=flush center,minimum height=1.55cm,anchor=west] (e21) at (6.10,-3.90) {Tempo di ritorno e orizzonte temporale\\[1pt]{\fontsize{7.6}{9}\selectfont\color{Muted}m14-ritorno}};
\draw[pEdge] (f2.east)--(e20.west);
\end{tikzpicture}
